%==============================================================================
%  (i)  Introduction
%==============================================================================
\section{Introduction}
\label{sec:intro}

% -----------------------------------------------------------------------------
% Move 1 -- the physical setting, and why anyone should care.
% Two or three paragraphs.  Keep it concrete: what the flow is, where it occurs,
% what quantity practice actually needs from it.
% -----------------------------------------------------------------------------
\TODO{Open on the configuration itself --- compound open-channel flow, the
main-channel/floodplain junction, the shear layer and the vertical-axis vortices
it sheds --- and on the engineering quantity that depends on them (conveyance,
lateral momentum and mass exchange, sediment and pollutant transport).
Cite the canonical experiments and the depth-averaged modelling tradition
\citep{Tominaga1991,ShionoKnight1991,Nezu2005}.}

% -----------------------------------------------------------------------------
% Move 2 -- what is already known.  Organise by *idea*, not by paper.
% One paragraph per strand of literature; each ends with what that strand still
% leaves open.  Do not write an annotated bibliography.
% -----------------------------------------------------------------------------
\TODO{Strand 1 --- experiments and eddy-resolving simulation of compound
channels: what has been measured and computed, and what is settled
\citep{Tominaga1991,Kara2012}.}

\TODO{Strand 2 --- coherent-structure decompositions of wall and free-surface
turbulence: POD, SPOD, and what they have revealed elsewhere
\citep{Lumley1970,Towne2018,SchmidtColonius2020}.}

\TODO{Strand 3 --- linear and resolvent analysis as a predictive rather than
descriptive tool: amplification of a turbulent mean flow, the non-normality
argument, and the cases where the leading resolvent mode has been shown to
recover the leading SPOD mode \citep{Trefethen1993,McKeonSharma2010,Towne2018}.}

% -----------------------------------------------------------------------------
% Move 3 -- the gap.  One paragraph, and it must be a *specific* gap, phrased so
% that the present work closes it.  "Has not been done" is weak; "the mechanism
% selecting X remains unexplained because Y has never been measured at the same
% (k_x, omega)" is strong.
% -----------------------------------------------------------------------------
\TODO{State the gap in one or two sentences.}

% -----------------------------------------------------------------------------
% Move 4 -- what this paper does, and the paper map.
% -----------------------------------------------------------------------------
\TODO{One paragraph: wall-resolved/wall-modelled LES of the compound channel at
$\Retau = \TODO{}$ and depth ratios $\Dr = \TODO{}$; SPOD of the resolved field
at each streamwise wavenumber; resolvent analysis of the same mean flow;
comparison of the two at matched $(\kx,\omega)$.}

The paper is organised as follows.
\Cref{sec:config} describes the flow configuration and the large-eddy
simulation.
\Cref{sec:spod} sets out the spectral proper orthogonal decomposition.
\Cref{sec:resolvent} develops the linear and resolvent analysis, including the
solver used to obtain the resolvent modes.
\Cref{sec:results} presents the numerical results, beginning with validation
against the reference data and continuing with \TODO{the analysis proper}.
\Cref{sec:conclusions} concludes.
